\documentclass[11pt]{article}
\usepackage[margin=1in]{geometry}
\usepackage[T1]{fontenc}
\usepackage{lmodern}
\usepackage{amsmath,booktabs,hyperref,microtype}
\hypersetup{hidelinks}
\makeatletter\let\tableinput\@@input\makeatother
% Generated; do not edit. Real corpus, partial identification.
\newcommand{\CipherCoreAssumedGuar}{no}
\newcommand{\CipherCoreAssumedPct}{5.3}
\newcommand{\CipherCoreRealGuar}{yes}
\newcommand{\CipherCoreRealPct}{3.9}
\newcommand{\CipherEhrShare}{14.8}
\newcommand{\CipherEmptyAssumedPct}{44.9}
\newcommand{\CipherEmptyRealPct}{33.0}
\newcommand{\CipherGamma}{1}
\newcommand{\CipherHighTotal}{128}
\newcommand{\CipherImgShare}{14.1}
\newcommand{\CipherLabShare}{18.8}
\newcommand{\CipherPharmaShare}{21.1}
\newcommand{\CipherRecords}{316}
\newcommand{\CipherResidualShare}{31.2}
\newcommand{\CipherTau}{7}
\newcommand{\CipherWeekTwoPlus}{30.7}


\title{Selection-robust clinical-harm bounds from a real ransomware-incident corpus,\\
and their effect on a control-adequacy certificate}
\author{}
\date{2026-09-20 (working note, version 0.1.0)}

\begin{document}
\maketitle

\begin{abstract}
Estimates of which hospital services ransomware harms most are usually read
straight off incident reports, ignoring that reports are a selective sample. We
take the CIPHER corpus (\CipherRecords{} coded patient-harm records from real
hospital-ransomware incidents) and, instead of treating its shares as unbiased,
compute \emph{partial-identification} bounds under bounded underreporting -- an
honest interval per clinical service rather than a point rate. Among the
\CipherHighTotal{} high-severity records (impact score $\ge \CipherTau{}$),
pharmacy (\CipherPharmaShare\%) and laboratory (\CipherLabShare\%) carry more
severe harm than the electronic health record (\CipherEhrShare\%), and
\CipherResidualShare\% of severe harm falls in domains the model does not
represent at all. Feeding these real-data-derived degradation intervals into the
clinical-impact certificate changes a conclusion: the hospital-core control bundle,
which is \emph{not} certified at the $5\%$ level under assumed inputs
(\CipherCoreAssumedPct\%), \emph{is} certified under the real-data intervals
(\CipherCoreRealPct\%). The contribution is the first inferential use of this
corpus and the selection-aware grounding, not a new theorem.
\end{abstract}

\section{Data and why bounds}
CIPHER (Straw, Tully, Dameff; CC-BY) codes patient-harm records from real
hospital-ransomware incident reports, each with a clinical-impact score, harm-onset
time point, and technical domain. It is a \textbf{convenience sample of reported
harms}: unreported harms are missing not at random, so a naive share is biased. We
therefore bound. Mapping technical domains to the four modelled clinical services
uses the repository's existing coverage map (Health Records~$\to$~EHR, Laboratory
Systems~$\to$~laboratory, ePrescribing~$\to$~pharmacy, Imaging~$\to$~imaging);
every other domain is kept as an explicit unmodelled residual, never dropped.

\section{Method}
Among high-severity records, let $n_s$ be service $s$'s count and $T$ the total
(including the residual). Under the assumption that any domain's true count is at
most $(1+\gamma)$ times its observed count, service $s$'s true share lies in
\[
\Big[\tfrac{n_s}{n_s+(1+\gamma)(T-n_s)},\ \tfrac{n_s(1+\gamma)}{n_s(1+\gamma)+(T-n_s)}\Big].
\]
At $\gamma=0$ this is the observed share; it widens monotonically with $\gamma$ and
stays in $[0,1]$. This is Manski-style partial identification; $\gamma$ is a
transparent worst-case dial, not a fitted selection model.

\section{Results}
Observed high-severity shares invert the usual intuition: pharmacy and laboratory
lead, not the EHR (Table~\ref{tab:shares}), and \CipherResidualShare\% of severe
harm is in domains outside the model (infrastructure, operating rooms,
communications). The harm-onset profile puts \CipherWeekTwoPlus\% of records at
Week~2 or later, corroborating that a short recovery horizon understates real harm.

Using the $\gamma=\CipherGamma$ intervals as the certificate's per-service
degradation changes a decision (Table~\ref{tab:compare}): undefended union clinical
outage moves from \CipherEmptyAssumedPct\% (assumed) to \CipherEmptyRealPct\%
(real-data, lower because roughly a third of severe harm is unmodelled), and the
hospital-core bundle crosses from \emph{not} certified at $5\%$
(\CipherCoreAssumedPct\%, guaranteed: \CipherCoreAssumedGuar) to certified
(\CipherCoreRealPct\%, guaranteed: \CipherCoreRealGuar). The point is not the
direction of the shift but that the certificate's answer is sensitive to this
input, and grounding it in real data --- with honest bounds --- is preferable to an
assumption.

\begin{table}[t]\centering
\caption{High-severity harm by clinical service: observed share and the
$\gamma=1$ partial-identification interval (percent of all high-severity records).}
\label{tab:shares}
\begin{tabular}{@{}lrrr@{}}\toprule
Service & Records & Observed share (\%) & $\gamma{=}1$ interval (\%)\\\midrule
\tableinput table_cipher_shares_rows.tex
\bottomrule\end{tabular}
\end{table}

\begin{table}[t]\centering
\caption{Clinical-impact certificate under assumed vs.\ CIPHER-derived degradation
intervals.}\label{tab:compare}
\begin{tabular}{@{}lrrl@{}}\toprule
Degradation input & Undefended union (\%) & Hospital-core union (\%) & Core guar.\ $\le$5\%\\\midrule
\tableinput table_cipher_compare_rows.tex
\bottomrule\end{tabular}
\end{table}

\section{Limitations}
CIPHER is reported harm, not incidence; the shares are relative
severity-attribution among reported harms, not probabilities from a random sample.
The single-$\gamma$ underreporting envelope is a worst-case device, not a fitted
model; counts are small, so intervals are wide -- the honest consequence. The
domain-to-service mapping is the repository's coverage mapping and leaves a large
unmodelled residual. This is a real-data grounding of one certificate input via
partial identification, not a new statistical method and not a validated estimate
of any hospital's risk.

\paragraph{Reproducibility.} CIPHER v1.0.1 is used unmodified (CC-BY), hashed in a
provenance manifest with every input; all numbers here are generated from the
committed tables.

\end{document}
